% controller.tex -- draft of the controller as flown (4 Oct 2026), input from Methodology.tex.
% Notation follows Source/src/model.tex: q_i points from the payload attachment to UAV i and
% T_i = lambda_i q_i. The code (mpc_planner, after Sun et al.) uses s_i = -q_i and t_i = lambda_i.
% Restructured 9 Oct (structure audit): the model moved to model.tex; planner, tracker and steady-error
% rejection are subsections of their own; the result numbers (the box's effect, K's measured values,
% the damping ratio) moved to Section 4.2; the thrust identification procedure moved here from
% hardware.tex; the cost weights moved to Appendix B. Notation: the horizon has N_h nodes over T_h
% (were N and T, which also named the team size and the tension).
% Its references (francis1976internal, pannocchia2003disturbance, pannocchia2007combined,
% maeder2009linear, verschueren2022acados) are in references.bib.

\subsubsection{Load--cable planner}
\label{sec:controller}

The controller is a two-level cascade of model predictive controllers (Section~\ref{sec:safety}; Figure~\ref{fig:owners}, Appendix~\ref{sec:app_method}).
A centralised planner on the model of Section~\ref{sec:model} plans the motion of the payload and of
every rod over a horizon. Each UAV's tracking MPC follows the planner's reference for that UAV,
accounting for the planned rod force, and sends a collective throttle and three body-rate commands
over the radio link to the flight controller, whose inner rate loop is not modified. Both levels are solved with acados
\cite{verschueren2022acados} in real-time iterations \cite{diehl2002rti}, warm-started from the previous solution,
with the QPs solved by HPIPM \cite{frison2020hpipm}.
Both horizons have $N_h = 20$ nodes over $T_h = 2$\,s. The team size selects the planner's solver and the
attachment geometry is a runtime parameter of it, so that a UAV leaving or joining switches the planner to a
pre-built problem of the new size (Section~\ref{sec:reconfig}). On the rig the planner starts no new iteration after 60\,ms of its
100\,ms tick, and a plan that has not converged by then is replaced by the previous one shifted by a tick (the measured
solve times are in Section~\ref{sec:res_compute}).

\paragraph{Optimal control problem.} At each planner tick the planner solves
\begin{equation}
\begin{aligned}
    \min_{\mathbf{x}_{0:N_h},\, \mathbf{u}_{0:N_h-1}} \quad & \sum_{k=0}^{N_h-1} \big\| \mathbf{y}(\mathbf{x}_k, \mathbf{u}_k) - \mathbf{y}_k^{\mathrm{ref}} \big\|_{W}^{2}
    + \big\| \mathbf{y}_{N_h}(\mathbf{x}_{N_h}) - \mathbf{y}_{N_h}^{\mathrm{ref}} \big\|_{W_{N_h}}^{2} \\
    \text{s.t.} \quad & \mathbf{x}_{k+1} = \mathbf{F}(\mathbf{x}_k, \mathbf{u}_k; \mathbf{d}), \\
    & C_0 \mathbf{x}_0 = C_0 \hat{\mathbf{x}}(t), \\
    & \lambda_{\min} \le \lambda_{i,k} \le \lambda_{\max}, \quad f_{\min} \le \|\mathbf{f}_{i,k}\| \le f_{\max}, \\
    & |\boldsymbol{\gamma}_{i,k}| \le \gamma_{\max}, \quad |\eta_{i,k}| \le \eta_{\max},
\end{aligned}
\label{eq:plan_ocp}
\end{equation}
with the output
$\mathbf{y} = [\mathbf{p}_L, \mathbf{v}_L, \mathbf{e}_R, \boldsymbol{\omega}_L, \{\mathbf{q}_i\}, \{\lambda_i\}, \{\mathbf{r}_i\}, \mathbf{u}]$,
where $\mathbf{e}_R = 2\,\mathrm{vec}(\mathbf{q}_L^{\mathrm{ref}} \otimes \mathbf{q}_L^{*})$ is the sign-invariant
attitude error. The thrust bounds are softened with slack. $\mathbf{F}$ is the integrated dynamics
over one node of $T_h/N_h = 0.1$\,s.

The initial condition pins only what motion capture measures ($C_0$ selects the payload pose and twist and each rod
direction); the rod rates and the tensions are free at node~0, warm-started from the previous solution.

\paragraph{References.} The payload reference $(\mathbf{p}_L^{\mathrm{ref}}, \mathbf{v}_L^{\mathrm{ref}},
\mathbf{a}_L^{\mathrm{ref}})$ comes from the lift schedule and the trajectory, and the rod references
follow from it by flatness. Holding the payload against gravity, its acceleration and the force
$\mathbf{d}$ needs the rods to support the effective gravity
\begin{equation}
    \mathbf{g}_{\mathrm{eff}} = g\mathbf{e}_3 + \mathbf{a}_{L,xy}^{\mathrm{ref}} - \mathbf{d}/m_L ,
\end{equation}
where $\mathbf{a}_{L,xy}^{\mathrm{ref}}$ is the horizontal part of the reference acceleration (the
trajectories are level). The nominal formation (each rod at the planned elevation, 45$^\circ$ until 7 October and 60$^\circ$ since, tensions
$\lambda_i^{0}$ balancing the payload weight) is rotated so that its vertical axis lies along
$\mathbf{g}_{\mathrm{eff}}$, and the tensions are scaled to $\lambda_i^{\mathrm{ref}} = \lambda_i^{0}\,
\|\mathbf{g}_{\mathrm{eff}}\| / g$. Of the weights (Table~\ref{tab:weights}, Appendix~\ref{sec:app_params}),
the one on the rod directions holds the formation's shape and the one on the rod rates damps rod swing.

The planner sends each tracker, at every node, its reference position and velocity from \eqref{eq:rod_kin}, the
planned thrust per unit mass from \eqref{eq:thrust_req} and the planned cable acceleration
$\mathbf{a}_{c,i} = -\lambda_i \mathbf{q}_i / m_q$, and the tracker interpolates the whole horizon at its own rate.

\subsection{Tracker, thrust map, measured pull and model integral}
\label{sec:tracker}

\subsubsection{UAV tracking MPC and thrust map}

Each UAV $i$ runs a tracking MPC on a 17-state model, in which the four stick channels sent to the
flight controller are themselves states and their rates of change are the inputs:
\begin{equation}
    \mathbf{x}_i = [\mathbf{p},\ \mathbf{q},\ \mathbf{v},\ \boldsymbol{\omega},\ \mathbf{c}], \qquad
    \mathbf{c} = [c_\phi, c_\theta, u_T, c_\psi], \qquad \mathbf{u}_i = \dot{\mathbf{c}},
\end{equation}
with
\begin{align}
    \dot{\mathbf{p}} &= \mathbf{v}, &
    \dot{\mathbf{q}} &= \tfrac{1}{2}\, \Omega(\boldsymbol{\omega})\, \mathbf{q}, \\
    \dot{\mathbf{v}} &= R(\mathbf{q})\, k_T u_T \mathbf{e}_3 - g\mathbf{e}_3 - c_z v_z \mathbf{e}_3 + \mathbf{a}_{c,i}, &
    \dot{\boldsymbol{\omega}} &= \frac{1}{\tau_\omega}\big(\boldsymbol{\omega}_{\mathrm{cmd}}(\mathbf{c}) - \boldsymbol{\omega}\big),
    \label{eq:trk_model}
\end{align}
where $\boldsymbol{\omega}_{\mathrm{cmd}}(\mathbf{c})$ is the flight controller's rate curve applied to the
roll, pitch and yaw sticks, and $\tau_\omega$ the time constant of its rate loop. Thrust is affine in
throttle: the throttle sent is $u_T + u_0(V)$, with the offset $u_0(V) = u_0 + \kappa (V_{\mathrm{ref}} - V)$
correcting for the battery voltage $V$ and $k_T$ the gain above it. The map was identified on the rig from free
hovers with hung masses and fitted against the mass carried and the voltage (Section~\ref{sec:thrust_id}).

The planned cable acceleration $\mathbf{a}_{c,i}$ enters the model as a known, time-varying parameter
along the horizon, so the tracker anticipates the rod's pull instead of reacting to it. The planned pull
is only as good as the planner's model of the rods. On the rig the rods flew at about 60$^\circ$ from
horizontal where the plan assumed 45$^\circ$, so each tracker now measures the pull and uses the plan only
for its shape along the horizon. The measured pull is the UAV's acceleration from motion capture (the
least-squares slope of its velocity over 80\,ms, low-pass filtered at 50\,ms) minus the modelled thrust,
gravity and drag, with the thrust passed through the same filter so that a throttle change does not read
as a change of pull. The horizon's planned pull is shifted by the difference between the measured and the
planned pull at its first node, capped at 12\,m/s$^2$, and the planner plans the rods at 60$^\circ$ to
match (effect on the rig: Section~\ref{sec:res_rig_integral}).

The tracker solves its own problem every 20\,ms over the same 2\,s horizon, with the throttle and the three stick
channels bounded and rate-limited and the attitude reference the direction of the planned thrust
(Appendix~\ref{sec:app_integral} gives it in full); only the first command of each solution is sent. Node~0 equals the
measured state, with the stick states taken from the commands actually sent; an earlier version relaxed it to the box
$\hat{\mathbf{x}}_i(1-\epsilon) \le \mathbf{x}_0 \le \hat{\mathbf{x}}_i(1+\epsilon)$, $\epsilon = 0.025$
(Section~\ref{sec:res_ic}).

\subsubsection{The model integral}
\label{sec:disturbance}

Any error that remains constant, from a model mismatch or an external force, is represented by the
constant force $\mathbf{d}$ on the payload in \eqref{eq:plan_trans}. $\mathbf{d}$ is a lumped disturbance:
for an external force (wind, a push, an unmodelled object in the basket) it is that force; for an
internal error (a thrust-map error, the rods' mass) it is the force that would have the same effect on
the payload, which need not exist physically. In the offset-free MPC literature this is an input
disturbance model \cite{pannocchia2003disturbance, maeder2009linear}.

\paragraph{Integral action through the model.} $\mathbf{d}$ is formed from the integral of the payload's
position error $\mathbf{e}_k = \mathbf{p}_L^{\mathrm{ref}} - \mathbf{p}_L$, updated at each planner tick
$k$ (period $T_s = 0.1$\,s):
\begin{equation}
    \mathbf{b}_{k+1} = \mathrm{sat}_{\bar{b}}\!\big(\mathbf{b}_k + k_i T_s\, G_k\, \mathbf{e}_k\big), \qquad
    \mathbf{d}_k = -K \mathbf{b}_k, \qquad K = \mathrm{diag}(K_{xy}, K_{xy}, K_z).
    \label{eq:lumped_int}
\end{equation}
Here $\mathrm{sat}_{\bar{b}}$ bounds each component to $\pm\bar{b}$, and the gate $G_k = \mathrm{diag}(g_{xy},
g_{xy}, g_z) \in \{0,1\}^3$ integrates only in steady, taut flight: lift complete, not descending or
landing, every rod taut, a fresh pose, the payload slow ($<0.05$\,m/s) and the miss under 0.25\,m. The
vertical component also integrates through a slow level orbit; the horizontal components require a
stationary reference, so a value learned in hover is held through a trajectory. A low payload ($e_z > 0$)
makes $d_z$ negative: the planner then believes the payload is heavier and plans more rod tension. Once
the payload rests on the floor at landing, $\mathbf{d}$ is set to zero. $K$ converts the integral from
metres to newtons and is set to the closed-loop stiffness of the cascade to a force on the payload,
measured from the steady displacement under a known force with the planned rod pull (Appendix~\ref{sec:app_params});
the integral's loop gain is then $k_i$, as in the reference-shift form below.

$\mathbf{d}_k$ is a parameter of the planner's problem, constant over the horizon and refreshed at every
tick. It enters the payload's translational dynamics \eqref{eq:plan_trans}, so the plan predicts the
payload's motion with the force included and chooses tensions and UAV positions that counter it. It also
enters the reference formation through $\mathbf{g}_{\mathrm{eff}}$, so the formation the cost draws the
solution towards is already tilted and tensioned to balance $\mathbf{d}$ and the cost does not oppose the
force balance.

\paragraph{What it does.} At an equilibrium with the gate open and no component of $\mathbf{b}$ saturated,
$\mathbf{b}_{k+1} = \mathbf{b}_k$ gives $\mathbf{e} = \mathbf{0}$ whatever the constant cause, by the internal model
principle \cite{francis1976internal}; with $K$ at the cascade's stiffness the believed force converges to the applied
one with time constant $1/k_i = 2.5$\,s, and when the force goes the payload overshoots by the displacement the force
caused at onset (Appendix~\ref{sec:app_integral}). With the measured rod pull, a constant vertical force is taken up at
each UAV within about a second, before the integral can act, so in the twin the integral is left with the horizontal
axes (Section~\ref{sec:res_force}). The bound $\bar{b} = 0.3$\,m, with a warning at 70\,\% of it, stops a wind-up, and a
force on the payload cannot represent a moment such as one UAV's thrust-map error, which that UAV's own tracker takes up.
An earlier form added the height integral to the height target instead; it flew on the rig on 1 and 8 October.
